\chapter{Conclusion}

MathDevMCP is proposed as a practical response to a recurring group-wide problem: mathematically intensive software and mathematically intensive documentation are developed together, but existing AI-assisted workflows do not reliably maintain the consistency, verification discipline, and long-range context needed for high-quality outcomes under tight deadlines. The group does not merely need better prose assistance. It needs a trustworthy environment for retrieval, verification, consistency checking, auditing, and controlled code generation.

The proposal recommends a disciplined path. Build a local Python toolkit that performs the real work of parsing, indexing, verification, and reporting. Expose a carefully selected subset of that functionality through an MCP server so Claude Code can use it interactively. Support the whole group through parallel CLI and batch-report interfaces. Evaluate every major capability using benchmarks that reflect the group's actual workflows: monograph support, code--document consistency, code generation correctness, documentation auditing, and usability under deadline pressure.

If executed in stages, MathDevMCP can become a durable internal capability rather than a one-off experiment. It can reduce duplicated review effort, lower the rate of unsupported or inconsistent mathematical claims, improve code--document traceability, and make long-form technical writing less brittle. For a group working on DSGE development, macrofinance, and similar mathematically structured projects, that would represent a meaningful improvement in both productivity and reliability.
